\chapter{基本群}
\label{chapter:ch09}

在本章中, 我们将用到前面各章中所学过的许多知识, 来构造代数拓扑学中的第一个工具——基本群. 我们用它来解决这样一些拓扑学问题: 给定两个空间, 什么时候它们之间存在同胚? 这个问题一般说来是困难的. 例如, 人们要花费很长时间才能看出为什么球面和环面不同胚, 或者为什么环面和双环面不同胚. 基本群就是解决这类问题的一种系统方法.

基本群的思想可以追溯到 Poincaré. 给定一个空间 $X$ 和一点 $x_0\in X$, 我们构造一个群 $\pi_1(X,x_0)$, 称为 $X$ 的以 $x_0$ 为基点的基本群. 这是一个从 $X$ 的拓扑构造出来的群. 如果 $X$ 和 $Y$ 是两个空间, 并且选定了基点 $x_0\in X$、$y_0\in Y$, 那么, 若存在一个同胚 $h\colon X\to Y$ 使得 $h(x_0)=y_0$, 则基本群 $\pi_1(X,x_0)$ 和 $\pi_1(Y,y_0)$ 便是同构的. 这样, 如果两个空间的基本群不同构, 我们便知道它们不同胚.

基本群的应用十分广泛. 我们用它来证明著名的代数基本定理, 证明 Brouwer 不动点定理, 证明 Borsuk-Ulam 定理, 比较紧致曲面的不同类型等. 我们还要研究覆叠空间, 它是将基本群的计算化为几何问题的一种有效手段.

基本群的构造涉及道路和道路同伦的概念, 这在第 51 节中给出. 基本群本身在第 52 节中定义. 覆叠空间在第 53 节中引入, 并在第 54 节中用来计算圆周的基本群. 从第 55 节到第 57 节, 我们给出基本群的若干经典应用. 第 58 节讨论伦型与同伦等价. 最后, 第 59 节证明 $S^n$ ($n\geqslant 2$) 是单连通的, 第 60 节计算环面、射影平面和双环面的基本群.

\subfile{Chapter-09/section-51}
\subfile{Chapter-09/section-52}
\subfile{Chapter-09/section-53}
\subfile{Chapter-09/section-54}
\subfile{Chapter-09/section-55}
\subfile{Chapter-09/section-56}
\subfile{Chapter-09/section-57}
\subfile{Chapter-09/section-58}
\subfile{Chapter-09/section-59}
\subfile{Chapter-09/section-60}
